\documentclass[11pt]{article}
\usepackage[T1]{fontenc}
\usepackage[utf8]{inputenc}
\usepackage{lmodern}
\usepackage[margin=1in]{geometry}
\usepackage{amsmath,amssymb,amsthm,mathtools}
\usepackage{booktabs,array,longtable}
\usepackage{microtype}
\usepackage{enumitem}
\usepackage{fancyhdr}
\usepackage{listings}
\usepackage[hidelinks]{hyperref}
\usepackage{bookmark}
\numberwithin{equation}{section}
\newtheorem{theorem}{Theorem}[section]
\newtheorem{proposition}[theorem]{Proposition}
\newtheorem{lemma}[theorem]{Lemma}
\newtheorem{corollary}[theorem]{Corollary}
\theoremstyle{definition}
\newtheorem{definition}[theorem]{Definition}
\newtheorem{conjecture}[theorem]{Conjecture}
\theoremstyle{remark}
\newtheorem{remark}[theorem]{Remark}
\DeclareMathOperator{\Res}{Res}
\DeclareMathOperator{\sgn}{sgn}
\DeclareMathOperator{\atanh}{atanh}
\newcommand{\fall}[2]{(#1)_{\underline{#2}}}
\newcommand{\coeff}[2]{[#1^{#2}]}
\newcommand{\Z}{\mathbb Z}
\newcommand{\Q}{\mathbb Q}
\newcommand{\R}{\mathbb R}
\newcommand{\U}{U}
\setlength{\headheight}{14pt}
\pagestyle{fancy}
\fancyhf{}
\fancyhead[L]{\small A290268: cancellations and nonvanishing}
\fancyhead[R]{\small Partial results}
\fancyfoot[C]{\thepage}
\setlength{\emergencystretch}{2em}
\lstset{basicstyle=\ttfamily\small,columns=fullflexible,breaklines=true,
  frame=single,keepspaces=true,showstringspaces=false}
\hypersetup{pdftitle={A290268: forced cancellations, nonvanishing regions, and certified computations},
  pdfauthor={Research note prepared for Vladimir Reshetnikov},
  pdfsubject={Rigorous partial results; the full conjecture is not proved}}

\title{The derivative-term conjecture for \texorpdfstring{$x^{x^2}$}{x raised to x squared}\\[4pt]
\large Forced cancellations, nonvanishing regions,\\and certified computations for OEIS A290268}
\author{Research note prepared for Vladimir Reshetnikov}
\date{September 20, 2026}

\begin{document}
\maketitle
\begin{abstract}
Let $a(n)$ be the number of nonzero monomials in the fully collected $n$th
 derivative of $x^{x^2}$. The conjecture in OEIS A290268 predicts
\[
 a(n)=\U(n):=\frac{n^2}{2}+n+1
 -(n\bmod2)\left(\frac12+\left\lfloor\frac{n+1}{8}\right\rfloor\right).
\]
This article does \emph{not} establish the equality for all $n$. It proves the
unconditional bound $a(n)\le\U(n)$ by identifying every cancellation needed to
explain the proposed formula. It also proves positivity throughout the upper
half of the coefficient triangle, classifies all zeros on the first two
logarithmic-deficit diagonals, and gives explicit quadratic lower bounds.
An exact generating-function identity and a finite-difference formula reduce
the full conjecture to a precise nonvanishing assertion. Integer arithmetic
verifies the support through $n=200$; two complete modular computations,
combined with the proved zero identities, certify $a(n)=\U(n)$ for every
$0\le n\le3000$. All source code, finite certificates, and reproduction
instructions accompany the article. The remaining nonvanishing step is
stated explicitly rather than assumed.
\end{abstract}

\section{The problem and the status of the result}\label{sec:introduction}
Put
\[
 f(x)=x^{x^2}=\exp(x^2\log x),\qquad x>0.
\]
The sequence A290268 counts the terms in the fully expanded derivative
$f^{(n)}(x)$, after identical monomials have been collected. The OEIS entry
was contributed by Vladimir Reshetnikov in October 2017; its formula section
labels both the generating function and the explicit expression as
conjectural. Peter Luschny's equivalent floor-function form is explicitly
conditional on that conjecture~\cite{oeis}.

For clarity, throughout this article the letter $\U$ denotes the
\emph{proposed} answer, not an already proved formula for $a$:
\begin{equation}\label{eq:U}
 \U(n)=\frac{n^2}{2}+n+1
 -(n\bmod2)\left(\frac12+\left\lfloor\frac{n+1}{8}\right\rfloor\right).
\end{equation}
Equivalently,
\begin{equation}\label{eq:U-parity}
 \U(2r)=2r^2+2r+1,\qquad
 \U(2r+1)=2(r+1)^2-\left\lfloor\frac{r+1}{4}\right\rfloor.
\end{equation}
The two central issues are different. One must explain why certain
coefficients vanish, and one must also exclude every other cancellation.
The first task is completed here. The second is completed in substantial
regions and in a large finite range, but not in full generality.

\begin{theorem}[Unconditional bounds]\label{thm:main-bounds}
For every integer $n\ge0$,
\begin{equation}\label{eq:bounds}
 \Lambda(n)\le a(n)\le\U(n),
 \qquad
 \Lambda(n)=
 \begin{cases}
 (3n^2+10n+8)/8,&n\text{ even},\\[2pt]
 (3n^2+12n+1)/8,&n\text{ odd}.
 \end{cases}
\end{equation}
In particular, $a(n)=\Theta(n^2)$. The proposed leading constant $1/2$ is
not established by these bounds.
\end{theorem}

The upper bound is proved in Section~\ref{sec:counting}; the lower bound is
proved in Section~\ref{sec:positivity}. Section~\ref{sec:computation} gives a
reproducible finite certification of the stronger equality through $3000$.
The exact outstanding assertion appears in Section~\ref{sec:remaining}.

\setcounter{tocdepth}{1}
\tableofcontents
\newpage

\section{A canonical polynomial expansion}\label{sec:canonical}
Set $z=x^2$ and $L=\log x$. Define polynomials $P_n\in\Z[z,L]$ by
\begin{equation}\label{eq:canonical}
 f^{(n)}(x)=f(x)x^{-n}P_n(x^2,\log x),\qquad P_0=1.
\end{equation}
Write
\begin{equation}\label{eq:c-def}
 P_n(z,L)=\sum_{k,j}c_{n,k,j}z^kL^j.
\end{equation}
Thus $c_{n,k,j}$ is the coefficient of
$x^{x^2+2k-n}(\log x)^j$ in the derivative. Counting terms means counting
nonzero $c_{n,k,j}$, rather than counting the summands generated by a
particular symbolic differentiation algorithm.

\begin{lemma}[Uniqueness of the collected expansion]\label{lem:unique}
A finite identity $\sum_{k,j}b_{k,j}x^{2k}(\log x)^j=0$ for all $x>0$
forces every $b_{k,j}$ to vanish. Consequently, the expansion
\eqref{eq:canonical}--\eqref{eq:c-def} is canonical.
\end{lemma}
\begin{proof}
Substitute $x=e^L$ and group the identity as
$\sum_k e^{2kL}B_k(L)=0$, with polynomial $B_k$. If $K$ is the largest
index with $B_K\ne0$, division by $e^{2KL}$ shows that $B_K(L)$ tends to
zero as $L\to+\infty$: all remaining terms are polynomials multiplied by
decaying exponentials. A nonzero polynomial cannot tend to zero at
$+\infty$. This contradiction eliminates $K$, and then all indices.
\end{proof}

The normalized derivative obeys a particularly small recurrence.
\begin{proposition}[Polynomial and coefficient recurrences]\label{prop:recurrence}
For every $n\ge0$,
\begin{equation}\label{eq:P-recurrence}
 P_{n+1}=\bigl(z(2L+1)-n\bigr)P_n
          +2z\frac{\partial P_n}{\partial z}
          +\frac{\partial P_n}{\partial L}.
\end{equation}
With nonexistent coefficients interpreted as zero,
\begin{equation}\label{eq:c-recurrence}
 \boxed{\;
 c_{n+1,k,j}=(2k-n)c_{n,k,j}+(j+1)c_{n,k,j+1}
             +c_{n,k-1,j}+2c_{n,k-1,j-1}.
 \;}
\end{equation}
For $n\ge1$, the only possible nonzero indices satisfy
\begin{equation}\label{eq:triangle}
 1\le k\le n,\qquad 0\le j\le k.
\end{equation}
\end{proposition}
\begin{proof}
We have $f'/f=x(2L+1)$, $dz/dx=2z/x$, and $dL/dx=1/x$.
Differentiate \eqref{eq:canonical} and factor out $f(x)x^{-n-1}$ to obtain
\eqref{eq:P-recurrence}. Taking the coefficient of $z^kL^j$ gives
\eqref{eq:c-recurrence}. Starting from $P_0=1$, the first derivative is
$P_1=z(2L+1)$. The four operations in \eqref{eq:P-recurrence} preserve
$k\ge1$ and $j\le k$, and increase $k$ by at most one. Induction proves
\eqref{eq:triangle} and integrality.
\end{proof}

For example,
\begin{align*}
 P_1&=z(2L+1),\\
 P_2&=z(2L+3)+z^2(2L+1)^2,\\
 P_3&=2z+z^2(12L^2+24L+9)+z^3(2L+1)^3.
\end{align*}
These have respectively $2$, $5$, and $8$ nonzero coefficients, in
agreement with the definition. Already $P_3$ illustrates a degree loss:
its $z$-coefficient has no $L$-term. Later cancellations are less obvious.
For instance,
\begin{equation}\label{eq:examples-cancel}
 [z^2]P_7(z,L)=28,\qquad [z^2]P_9(z,L)=-864L.
\end{equation}
The missing $L$-term in the first polynomial and the missing constant in
the second will follow from the same reflection principle.

\section{Exact formulas for individual coefficients}\label{sec:formulas}
\subsection{An exponential generating function}
Define two formal power series
\begin{equation}\label{eq:AB}
 A(t)=t(2+t)=(1+t)^2-1,\qquad
 B(t)=(1+t)^2\log(1+t).
\end{equation}
Formal coefficient extraction is denoted by $[t^n]$. All logarithms in
formal series are expanded at the origin.

\begin{proposition}[Generating function]\label{prop:egf}
The complete coefficient array is described by
\begin{equation}\label{eq:egf}
 \sum_{n\ge0}P_n(z,L)\frac{t^n}{n!}
   =\exp\bigl(z(A(t)L+B(t))\bigr).
\end{equation}
In particular, if $m=k-j\ge0$, then
\begin{equation}\label{eq:coefficient-AB}
 \boxed{\quad
 c_{n,k,j}=\frac{n!}{j!m!}[t^n]A(t)^j B(t)^m.
 \quad}
\end{equation}
\end{proposition}
\begin{proof}
Taylor expansion at the point $x$ gives
\[
 \frac{f(x(1+t))}{f(x)}
 =\sum_{n\ge0}\frac{x^nf^{(n)}(x)}{f(x)}\frac{t^n}{n!}.
\]
The logarithm of the ratio on the left is
\[
 x^2\bigl(((1+t)^2-1)\log x+(1+t)^2\log(1+t)\bigr).
\]
This proves \eqref{eq:egf} after substitution of $z,L$. Alternatively,
its coefficients satisfy \eqref{eq:P-recurrence}, so the identity also
holds with $z,L$ algebraically independent. Expanding separately the two
exponentials $\exp(zA(t)L)$ and $\exp(zB(t))$ gives
\eqref{eq:coefficient-AB}.
\end{proof}

Because both $A$ and $B$ start with a nonzero term of degree one, their
product in \eqref{eq:coefficient-AB} starts at degree $k$. This independently
explains $k\le n$. The difficulty is that $B$ does not have only positive
coefficients:
\[
 B(t)=t+\frac32t^2+\frac13t^3-\frac1{12}t^4+\frac1{30}t^5-\cdots.
\]
Consequently, the coefficient formula is not itself a proof of
nonvanishing.

\subsection{Falling factorials and finite differences}
Let
\[
 F_n(u)=\fall{u}{n}=u(u-1)\cdots(u-n+1),\qquad F_0(u)=1,
\]
and let $\Delta_2G(u)=G(u+2)-G(u)$. Set
$Q_{n,j}(u)=\Delta_2^jF_n(u)$.

\begin{proposition}[Finite-difference formula]\label{prop:difference}
For $m=k-j\ge0$,
\begin{align}
 c_{n,k,j}
 &=\frac{1}{j!m!}
       \left.\frac{d^m}{du^m}Q_{n,j}(u)\right|_{u=2m}
       \label{eq:difference}\\
 &=\frac1{j!}[v^m]\sum_{h=0}^j(-1)^{j-h}\binom jh
       F_n(2m+2h+v).\label{eq:difference-coeff}
\end{align}
\end{proposition}
\begin{proof}
The generalized binomial theorem gives
$n![t^n](1+t)^u=F_n(u)$. Differentiating $m$ times in $u$ introduces
$\log^m(1+t)$. Taking $j$ forward differences of step two introduces
$((1+t)^2-1)^j=A(t)^j$. Evaluation at $u=2m$ introduces
$(1+t)^{2m}$, which combines with the logarithm to give $B(t)^m$.
Equation~\eqref{eq:coefficient-AB} now proves \eqref{eq:difference};
Taylor's formula proves \eqref{eq:difference-coeff}.
\end{proof}

Derivatives of falling factorials also occur in the theory of shifted
Stirling numbers and elementary symmetric harmonic sums; Poon gives a
broader account of these connections~\cite{poon}. No general nonvanishing
theorem from that literature is being assumed here. Formula
\eqref{eq:difference-coeff} is also an independent way to compute a
coefficient using integer polynomial products.

\section{A transformation proving positivity}\label{sec:positivity}
A fractional-linear change of variable turns a large region of the
coefficient array into a manifestly positive expression. Define
\begin{equation}\label{eq:H}
 H(u)=\frac{\atanh u}{u}=\sum_{r\ge0}\frac{u^{2r}}{2r+1},
 \qquad H(0)=1.
\end{equation}
The value at zero is understood by its formal power series.

\begin{proposition}[Transformed coefficient formula]\label{prop:mobius}
For $m=k-j\ge0$ and $n\ge k$,
\begin{equation}\label{eq:mobius}
 c_{n,k,j}=\frac{n!}{j!m!}\,2^{\,2j+m-n}
 [u^{n-k}](1-u)^{n-2k-1}(1+u)^{2m}H(u)^m.
\end{equation}
Negative integer powers of $1-u$ have their usual formal expansions.
\end{proposition}
\begin{proof}
In the residue expression for $[t^n]A(t)^jB(t)^m$, substitute
$t=2u/(1-u)$. Then
\[
 1+t=\frac{1+u}{1-u},\quad
 A(t)=\frac{4u}{(1-u)^2},\quad
 B(t)=\frac{2(1+u)^2}{(1-u)^2}\atanh u,
 \quad dt=\frac{2\,du}{(1-u)^2}.
\]
Using $[t^n]G(t)=\Res_{t=0}G(t)t^{-n-1}\,dt$, the constant factor becomes
$2^{2j+m-n}$. The remaining power of $1-u$ is $n-2k-1$.
The factors $u^j(\atanh u)^m$ contribute $u^kH(u)^m$, so the resulting
residue extracts the coefficient of $u^{n-k}$. Multiplying by
$n!/(j!m!)$ proves the formula. The formal change-of-variable rule applies
because $2u/(1-u)$ has zero constant term and nonzero linear coefficient.
\end{proof}

\begin{theorem}[A nonvanishing region]\label{thm:positive}
For $1\le k\le n$ and $0\le j\le k$,
\begin{equation}\label{eq:positive}
 \begin{split}
 n\le2k&\quad\Longrightarrow\quad c_{n,k,j}>0,\\
 n=2k+1,\ j<k&\quad\Longrightarrow\quad c_{n,k,j}>0.
 \end{split}
\end{equation}
\end{theorem}
\begin{proof}
If $n\le2k$, the factor $(1-u)^{n-2k-1}$ is $(1-u)^{-q}$ with $q\ge1$.
Every coefficient of this factor is strictly positive. The other two
factors in \eqref{eq:mobius} have nonnegative coefficients and constant
term one. Their product therefore has a strictly positive coefficient in
every nonnegative degree.

If $n=2k+1$, that factor disappears. The condition $j<k$ means $m\ge1$.
The series $H(u)^m$ has a strictly positive coefficient in every even
degree. The polynomial $(1+u)^{2m}$ has positive constant and linear
coefficients. An even target degree can be supplied entirely by $H(u)^m$;
an odd target degree can use the linear term of $(1+u)^{2m}$ and an even
degree from $H(u)^m$. Thus the target coefficient of degree $n-k$ is
strictly positive in both cases. The prefactor in \eqref{eq:mobius} is
positive.
\end{proof}

\begin{corollary}[Quadratic lower bounds]\label{cor:lower}
The lower bound in Theorem~\ref{thm:main-bounds} holds.
\end{corollary}
\begin{proof}
For $n=2r\ge2$, Theorem~\ref{thm:positive} supplies all coefficients with
$r\le k\le2r$, giving
\[
 a(2r)\ge\sum_{k=r}^{2r}(k+1)=\frac{(r+1)(3r+2)}2.
\]
For $n=2r+1$ with $r\ge1$, it supplies all coefficients with
$r+1\le k\le2r+1$ and the $r$ coefficients $0\le j<r$ at $k=r$.
Hence
\[
 a(2r+1)\ge r+\sum_{k=r+1}^{2r+1}(k+1)
 =\frac{3(r+1)(r+2)}2-1.
\]
These are precisely \eqref{eq:bounds}. The cases $n=0,1$ follow directly
from $P_0=1$ and $P_1=z(2L+1)$.
\end{proof}

\section{The two forced families of zeros}\label{sec:zeros}
\subsection{Degree-forced zeros on the top diagonal}
Setting $j=k$, and therefore $m=0$, in \eqref{eq:coefficient-AB} yields
an explicit coefficient.
\begin{proposition}[Top logarithmic degree]\label{prop:top}
For $k\ge1$ and $k\le n\le2k$,
\begin{equation}\label{eq:top}
 c_{n,k,k}=\frac{n!}{k!}\binom{k}{n-k}2^{2k-n}>0.
\end{equation}
For $n>2k$, one has $c_{n,k,k}=0$.
\end{proposition}
\begin{proof}
The relevant series is the polynomial
$A(t)^k=t^k(2+t)^k$, of degree $2k$. Its coefficient of $t^n$ is
$\binom{k}{n-k}2^{2k-n}$ in the indicated range and zero above it.
\end{proof}
For a fixed $n\ge1$, this removes exactly
$\lfloor(n-1)/2\rfloor$ cells from the possible triangle.

\subsection{Reflection-forced cancellations}
The less obvious zeros follow from symmetry about the midpoint of the
roots of a falling factorial.
\begin{lemma}[Reflection of a finite difference]\label{lem:reflection}
For all $n,j\ge0$,
\begin{equation}\label{eq:reflection}
 Q_{n,j}(n-2j-1-u)=(-1)^{n-j}Q_{n,j}(u).
\end{equation}
\end{lemma}
\begin{proof}
First,
$F_n(n-1-u)=(-1)^nF_n(u)$, by reversing the order of the factors.
Expand the left side of \eqref{eq:reflection}:
\[
 \sum_{h=0}^j(-1)^{j-h}\binom jh
 F_n(n-2j-1-u+2h).
\]
Reflection turns its last factor into
$(-1)^nF_n(u+2(j-h))$. Reindexing by $\ell=j-h$ leaves the factor
$(-1)^{n-j}$ times the usual expansion of $\Delta_2^jF_n(u)$.
\end{proof}

\begin{theorem}[Reflection zeros]\label{thm:reflection-zeros}
Suppose $1\le k\le n$ and $0\le j<k$. Then
\begin{equation}\label{eq:reflection-zeros}
 k\text{ even},\qquad n=4k-2j+1
 \quad\Longrightarrow\quad c_{n,k,j}=0.
\end{equation}
\end{theorem}
\begin{proof}
Put $m=k-j$. The reflection center in Lemma~\ref{lem:reflection} is
$(n-2j-1)/2$. Under the displayed condition this equals $2m$, precisely
the evaluation point in \eqref{eq:difference}. The parity of
$Q_{n,j}(2m+v)$ as a polynomial in $v$ is $n-j$. Its coefficient of $v^m$
therefore vanishes whenever $n-j-m=n-k$ is odd. The relation
$n=4k-2j+1$ makes $n$ odd, so $n-k$ is odd when $k$ is even.
Equation~\eqref{eq:difference-coeff} completes the proof.
\end{proof}

There is also a direct proof using \eqref{eq:mobius}. The relation in
\eqref{eq:reflection-zeros} gives $n-2k-1=2m$, so the series whose
coefficient is extracted becomes
\begin{equation}\label{eq:even-kernel}
 (1-u^2)^{2m}H(u)^m.
\end{equation}
This series is even. When $k$ is even, $n-k$ is odd, and its coefficient
of that degree is zero. This second proof makes the cancellation visible
without evaluating any large integer coefficient.

\begin{remark}[What symmetry does not say]\label{rem:symmetry-gap}
The theorem is an implication, not an equivalence. An even power series
has zero odd coefficients, but it can also have zero even coefficients.
Similarly, away from the reflection center, a derivative of a polynomial
can vanish for reasons unrelated to symmetry. Excluding such additional
zeros is a separate requirement of the full conjecture.
\end{remark}

\section{Counting the forced zeros}\label{sec:counting}
For $n\ge1$, the triangle \eqref{eq:triangle} has
\begin{equation}\label{eq:triangle-size}
 T(n)=\sum_{k=1}^n(k+1)=\frac{n(n+3)}2
\end{equation}
cells. Let $Z_n^{\mathrm{top}}$ be the top-diagonal cells removed by
Proposition~\ref{prop:top}, and let $Z_n^{\mathrm{ref}}$ be the cells in
Theorem~\ref{thm:reflection-zeros}. These sets are disjoint, since the
second requires $j<k$.

\begin{lemma}[Number of reflection zeros]\label{lem:reflection-count}
For even $n$, $Z_n^{\mathrm{ref}}$ is empty. For odd $n$,
\begin{equation}\label{eq:reflection-count}
 |Z_n^{\mathrm{ref}}|=\left\lfloor\frac{n+1}{8}\right\rfloor.
\end{equation}
\end{lemma}
\begin{proof}
Reflection zeros have odd $n$. Write $n=2r+1$ and $k=2\ell$.
The reflection relation determines
$j=2k-r=4\ell-r$. The inequalities $0\le j<k$ become
\[
 \frac r4\le\ell<\frac r2.
\]
For $r\ge1$, these automatically give $\ell\ge1$ whenever the interval
contains an integer. Its number of integers is
\[
 \left\lceil\frac r2\right\rceil-
 \left\lceil\frac r4\right\rceil
 =\left\lfloor\frac{r+1}{4}\right\rfloor,
\]
as is checked in the four residue classes of $r$ modulo four. For $r=0$
both sides are zero. Since $(r+1)/4=(n+1)/8$, the result follows.
\end{proof}

\begin{theorem}[The conjectured expression is an upper bound]\label{thm:upper}
For every $n\ge0$,
\begin{equation}\label{eq:upper}
 a(n)\le\U(n).
\end{equation}
Equality holds for a given $n$ if and only if the coefficient triangle has
no zeros other than $Z_n^{\mathrm{top}}\cup Z_n^{\mathrm{ref}}$.
\end{theorem}
\begin{proof}
For even $n=2r\ge2$, subtraction of the top-diagonal zeros from
\eqref{eq:triangle-size} gives
\[
 2r^2+3r-(r-1)=2r^2+2r+1.
\]
For odd $n=2r+1$, subtraction of both disjoint families gives
\[
 2r^2+5r+2-r-\left\lfloor\frac{r+1}{4}\right\rfloor
 =2(r+1)^2-\left\lfloor\frac{r+1}{4}\right\rfloor.
\]
These are \eqref{eq:U-parity}. Additional zeros could only decrease the
number of nonzero coefficients, proving the inequality and the stated
equivalence. Finally $a(0)=\U(0)=1$.
\end{proof}

For example, at $n=7$ there are $35$ potential cells, three top-diagonal
zeros, and the reflection zero $(k,j)=(2,1)$, leaving $31$. At $n=9$,
there are $54$ potential cells, four top-diagonal zeros, and the reflection
zero $(2,0)$, leaving $49$. The exact expansions in
\eqref{eq:examples-cancel} display those two different cancellations.
The floor term in the proposed formula is thus explained by a proved
symmetry, not merely by fitting initial values.

\section{Nonvanishing on the first two deficit diagonals}\label{sec:small-m}
Call $m=k-j$ the \emph{logarithmic deficit}. This section gives a complete
zero classification for $m=1$ and $m=2$. Together with
Theorem~\ref{thm:positive}, it narrows the unproved portion of the problem.

\subsection{A beta-integral representation}
Suppose $n>2k$, $m=k-j\ge1$, and put
\begin{equation}\label{eq:d}
 d=n-2k-1\ge0.
\end{equation}
For sufficiently small complex $\varepsilon$, define
\begin{equation}\label{eq:I}
 I_{m,j,d}(\varepsilon)=
 \int_0^1 t^{2m+\varepsilon}(1-t)^{d-\varepsilon}(2t-1)^j\,dt.
\end{equation}
The integral and all derivatives required below converge uniformly on
compact subsets of $|\varepsilon|<1/2$.

\begin{lemma}[Integral formula]\label{lem:beta}
Under these assumptions,
\begin{equation}\label{eq:beta-Q}
 Q_{n,j}(2m+\varepsilon)
 =(-1)^{n-1}n!\frac{\sin(\pi\varepsilon)}{\pi}
    I_{m,j,d}(\varepsilon),
\end{equation}
and consequently
\begin{equation}\label{eq:beta-c}
 c_{n,k,j}=\frac{(-1)^{n-1}n!}{j!}
 [\varepsilon^m]\left(
 \frac{\sin(\pi\varepsilon)}{\pi}I_{m,j,d}(\varepsilon)\right).
\end{equation}
\end{lemma}
\begin{proof}
Euler's reflection formula and beta integral~\cite{dlmf-reflection,dlmf-beta}
give, initially away from integer $\alpha$,
\begin{align*}
 F_n(\alpha)
 &=(-1)^{n-1}\frac{\sin(\pi\alpha)}{\pi}
      \Gamma(\alpha+1)\Gamma(n-\alpha)\\
 &=(-1)^{n-1}n!\frac{\sin(\pi\alpha)}{\pi}
      \int_0^1t^\alpha(1-t)^{n-\alpha-1}\,dt.
\end{align*}
In the $j$th difference put $\alpha=2m+2h+\varepsilon$. The sine factor
is independent of the integer $h$. The finite sum inside the integral is
\[
 \sum_{h=0}^j(-1)^{j-h}\binom jh
     \left(\frac{t}{1-t}\right)^{2h}
 =\frac{(2t-1)^j}{(1-t)^{2j}}.
\]
The remaining exponent of $1-t$ is $d-\varepsilon$, proving
\eqref{eq:beta-Q}. All terms are analytic near $\varepsilon=0$, so the
identity extends to zero as well. Taking the coefficient of
$\varepsilon^m$ and using \eqref{eq:difference-coeff} gives
\eqref{eq:beta-c}.
\end{proof}

\subsection{Deficit one}
\begin{theorem}[Complete signs for $m=1$]\label{thm:m1}
Suppose $m=k-j=1$ and $n>2k$, and define $d$ by \eqref{eq:d}. Then
\begin{equation}\label{eq:m1-sign}
 \sgn c_{n,k,j}=(-1)^{n-1}
 \begin{cases}
 1,&j\text{ even},\\
 \sgn(2-d),&j\text{ odd}.
 \end{cases}
\end{equation}
In particular the only zeros in this range have $j$ odd and $d=2$.
\end{theorem}
\begin{proof}
The coefficient of $\varepsilon$ in \eqref{eq:beta-c} is
\begin{equation}\label{eq:m1-integral}
 c_{n,k,j}=\frac{(-1)^{n-1}n!}{j!}
 \int_0^1t^2(1-t)^d(2t-1)^j\,dt.
\end{equation}
For even $j$, the integrand is nonnegative and positive except at finitely
many points, so the integral is strictly positive. For odd $j$, pair $t$
with $1-t$ and integrate over $1/2<t<1$. The paired integrand is
\[
 (2t-1)^j\bigl(t^2(1-t)^d-t^d(1-t)^2\bigr).
\]
The ratio of the two positive terms in parentheses is
$(t/(1-t))^{2-d}$. Since $t/(1-t)>1$, their difference is identically
zero when $d=2$ and otherwise has the strict sign of $2-d$ throughout
the open interval. This proves the sign formula.
\end{proof}
These zeros are exactly the previously identified reflection zeros:
$j$ odd means $k=j+1$ is even, and $d=2$ means
$n=2k+3=4k-2j+1$.

\subsection{Deficit two}
\begin{theorem}[Complete signs for $m=2$]\label{thm:m2}
Suppose $m=k-j=2$ and $n>2k$. Then
\begin{equation}\label{eq:m2-sign}
 \sgn c_{n,k,j}=(-1)^{n-1}
 \begin{cases}
 1,&j\text{ odd},\\
 \sgn(4-d),&j\text{ even}.
 \end{cases}
\end{equation}
The only zeros in this range have $j$ even and $d=4$.
\end{theorem}
\begin{proof}
The coefficient of $\varepsilon^2$ in \eqref{eq:beta-c} is
\begin{equation}\label{eq:m2-integral}
 c_{n,k,j}=\frac{(-1)^{n-1}n!}{j!}
 \int_0^1t^4(1-t)^d(2t-1)^j
       \log\frac{t}{1-t}\,dt.
\end{equation}
For odd $j$, the factors $(2t-1)^j$ and $\log(t/(1-t))$ have the same
sign, so the integral is strictly positive. For even $j$, their product
is antisymmetric under $t\mapsto1-t$. Pairing the integral as before
leaves, on $1/2<t<1$,
\[
 (2t-1)^j\log\frac{t}{1-t}
 \bigl(t^4(1-t)^d-t^d(1-t)^4\bigr).
\]
The first two factors are positive there. The bracket has the strict
sign of $4-d$, or is identically zero when $d=4$. This proves the claim.
\end{proof}
Again the zero condition is exactly reflection: $k=j+2$ is even and
$n=2k+5=4k-2j+1$.

For $n\le2k$, positivity has already been proved, and the borderline
$n=2k+1$ is included in both sign arguments. Thus these two theorems,
together with Proposition~\ref{prop:top}, classify all zeros for
$k-j\le2$ throughout the possible triangle.

\subsection{Why the same sign proof does not settle higher deficits}
At $m=3$, formula \eqref{eq:beta-c} gives the integral of
\begin{equation}\label{eq:m3-kernel}
 t^6(1-t)^d(2t-1)^j
 \left(\frac12\log^2\frac{t}{1-t}-\frac{\pi^2}{6}\right)
\end{equation}
multiplied by $(-1)^{n-1}n!/j!$. The last factor changes sign inside
$(0,1)$. Thus positivity of the weight, or pairing $t$ with $1-t$, no
longer excludes cancellation. This is a concrete obstruction to simply
extending the proofs above. It is not a counterexample to the conjecture.

\section{Generating functions and equivalent formulas for the upper bound}
\label{sec:gf}
All identities in this section are proved identities for $\U(n)$.
They become identities for $a(n)$ only after the remaining nonvanishing
assertion is supplied, or when restricted to a certified finite range.

\subsection{Derivation of the rational generating function}
By \eqref{eq:U-parity}, the even and odd subsequences have generating
functions
\begin{align}
 E(y)&=\sum_{r\ge0}\U(2r)y^r
       =\frac{(1+y)^2}{(1-y)^3},\label{eq:E}\\
 O(y)&=\sum_{r\ge0}\U(2r+1)y^r
       =\frac{2(1+y)}{(1-y)^3}
          -\frac{y^3}{(1-y)(1-y^4)}.\label{eq:O}
\end{align}
For the last term, write
$\lfloor(r+1)/4\rfloor=\sum_{h\ge1}\mathbf1_{r\ge4h-1}$ and sum each
geometric tail. Therefore
\begin{align}
 G(t):=\sum_{n\ge0}\U(n)t^n
 &=E(t^2)+tO(t^2)\notag\\
 &=\frac{1+t^2}{(1-t)^3(1+t)}
       -\frac{t^7}{(1-t^2)(1-t^8)}\notag\\
 &=\frac{N(t)}{(1-t)(1-t^2)(1-t^8)},\label{eq:GF}
\end{align}
where
\begin{equation}\label{eq:N}
 N(t)=1+t+2t^2+2t^3+2t^4+2t^5+2t^6+t^7+2t^8+t^9.
\end{equation}
Indeed,
$N(t)=(1+t^2)(1+t+\cdots+t^7)-t^7(1-t)$.
This is exactly the rational function proposed in the OEIS entry~\cite{oeis}.
The derivation shows why period eight appears: it counts the even values
of $k$ on the reflection line, in addition to the parity split of the
ordinary triangle.

\subsection{The eight constituent polynomials}
Writing $n=8q+s$ with $0\le s\le7$, \eqref{eq:U} becomes the following
degree-two quasipolynomial.
\begin{center}
\begin{tabular}{cl}
\toprule
$s$ & $\U(8q+s)$\\
\midrule
0 & $32q^2+8q+1$\\
1 & $32q^2+15q+2$\\
2 & $32q^2+24q+5$\\
3 & $32q^2+31q+8$\\
4 & $32q^2+40q+13$\\
5 & $32q^2+47q+18$\\
6 & $32q^2+56q+25$\\
7 & $32q^2+63q+31$\\
\bottomrule
\end{tabular}
\end{center}
The trigonometric expression from the OEIS entry is another way to write
the same eight polynomials:
\begin{equation}\label{eq:trig}
 \U(n)=\frac{8n^2+15n+14+(n+2)(-1)^n
 +(2-4\sqrt2\sin(\pi n/4))\sin(\pi n/2)}{16}.
\end{equation}
To verify the equivalence without an approximation, substitute $n=8q+s$.
The sine factors take their exact values at multiples of $\pi/4$ and
$\pi/2$. Substitution in each of the eight residue classes yields the
polynomials in the table. In particular, no irrational value remains.

\subsection{Linear recurrences}
Multiplication of \eqref{eq:GF} by its denominator yields
\begin{equation}\label{eq:rec11}
 \begin{split}
 \U(n)={}&\U(n-1)+\U(n-2)-\U(n-3)+\U(n-8)\\
         &-\U(n-9)-\U(n-10)+\U(n-11),\qquad n\ge11.
 \end{split}
\end{equation}
Multiplying the denominator by $1+t$ gives the convenient order-twelve
form recorded in the OEIS entry:
\begin{equation}\label{eq:rec12}
 \U(n)=2\U(n-2)-\U(n-4)+\U(n-8)-2\U(n-10)+\U(n-12),
 \qquad n\ge12.
\end{equation}
The twelve initial values are
\[
 1,2,5,8,13,18,25,31,41,49,61,71.
\]
Using these recurrences to \emph{generate} the claimed values of $a(n)$
would be circular in a proof of the conjecture. The computations in the
next section instead use the independently derived coefficient recurrence
\eqref{eq:c-recurrence}.

\section{Exact finite verification and certificates}\label{sec:computation}
\subsection{Integer arithmetic and independent formula checks}
The accompanying Python program applies \eqref{eq:c-recurrence} with
arbitrary-precision integers. Through $n=200$, it checks not only the
term count but the entire zero pattern: a coefficient is zero if and only
if it belongs to one of the two proved zero families. This examines
$1{,}373{,}500$ cells with $1\le k\le n\le200$, $0\le j\le k$.

For $0\le n\le16$, the program independently recomputes $969$ triples
$(n,k,j)$, including the $k=0$ boundary, using each of
\eqref{eq:coefficient-AB}, \eqref{eq:difference-coeff}, and
\eqref{eq:mobius}. Fractions are exact rational numbers. All three formulas
agree with the integer recurrence. The positivity and deficit-one and
-deficit-two sign theorems are also checked throughout the integer range.
These checks test the derivations and their implementation; the proofs in
the preceding sections do not depend on them.

\subsection{Why modular arithmetic gives a rigorous finite certificate}
A nonzero residue modulo any integer $M\ge2$ proves that the original
integer is nonzero. A zero residue, in contrast, might mean only that
$M$ divides a nonzero integer. We use two moduli,
\begin{equation}\label{eq:moduli}
 M_1=1{,}000{,}000{,}007,\qquad M_2=1{,}000{,}000{,}009.
\end{equation}
Their primality is not needed for this argument: only exact reduction of
an integer recurrence is used, with no modular division.

For each modulus, a C++ program computes every row through $n=3000$ and
records every zero residue outside the two mathematically proved zero
families. There are three such residues modulo $M_1$ and four modulo
$M_2$. The two lists are disjoint. Their union and the corresponding
residues are as follows.
\begin{center}
\small
\begin{tabular}{rrr|rr}
\toprule
$n$ & $k$ & $j$ & $c_{n,k,j}\bmod M_1$ & $c_{n,k,j}\bmod M_2$\\
\midrule
1471&1046&521&847578491&0\\
1701&610&505&0&439228697\\
1854&1713&549&0&366891231\\
2336&1689&145&0&782191592\\
2543&1785&663&405694396&0\\
2807&2566&502&607368440&0\\
2842&2673&2118&474540977&0\\
\bottomrule
\end{tabular}
\end{center}
Every entry in this table is therefore a nonzero integer coefficient.
All other cells outside the proved zero families have a nonzero residue
in both runs. Each modulus run checks $4{,}509{,}002{,}501$ potential
cells, including the single cell at $n=0$.

\begin{proposition}[Finite computational certification]\label{prop:finite}
The accompanying exact computations, together with
Proposition~\ref{prop:top} and Theorem~\ref{thm:reflection-zeros}, certify
\begin{equation}\label{eq:finite}
 a(n)=\U(n)\qquad\text{for every integer }0\le n\le3000.
\end{equation}
\end{proposition}
\begin{proof}[Certificate argument]
The two analytic results prove that every predicted zero is an actual
integer zero. For any other coefficient in the finite range, membership
in neither, one, or both modular zero lists exhausts the possibilities.
The intersection of the lists is empty. Thus at least one of its two
residues is nonzero, proving that the integer coefficient is nonzero.
Theorem~\ref{thm:upper} now gives the equality for every row in the range.
\end{proof}
This is a deterministic finite verification, not a probabilistic claim
based on the size of the moduli. It is also not a formal proof-assistant
certificate: it relies on the supplied ordinary programs and their
execution, which can be independently reproduced.

\subsection{Implementation details and limitations}
The C++ program stores two coefficient rows in arrays of unsigned
32-bit integers. An update is evaluated in signed 64-bit arithmetic
before reduction. Its explicit input bounds guarantee that these
intermediates cannot overflow. It never approximates logarithms or
transcendental values. The arrays have $O(N^2)$ entries and the computation
uses $O(N^3)$ fixed-modulus arithmetic operations through order $N$.

The Python implementation uses a different update organization: it
scatters each existing coefficient into its four successors, whereas
the C++ implementation gathers four contributions into each destination.
The small-range generating-function and finite-difference calculations
provide further independent comparisons. The archive retains both raw
modular zero lists, row-by-row support statistics, explicit nonzero
witnesses, the combined certificate, and the exact small-range counts.

No finite range, regardless of its size, proves the nonvanishing assertion
for all $n$. In particular, the certificate supplies neither an induction
step beyond $3000$ nor a theorem bounding where a first additional zero
could occur.

\section{The precise missing step}\label{sec:remaining}
The coefficient problem can now be stated independently of the original
transcendental function.
\begin{conjecture}[Complete support classification]\label{conj:support}
Let $n\ge1$, $1\le k\le n$, $0\le j\le k$, and $m=k-j$. Then
\begin{equation}\label{eq:classification}
 \frac{1}{j!m!}
 \left.\frac{d^m}{du^m}\Delta_2^j\fall{u}{n}\right|_{u=2m}=0
\end{equation}
if and only if at least one of the following holds:
\begin{equation}\label{eq:only-zeros}
 \bigl(j=k\text{ and }n>2k\bigr),
 \qquad
 \bigl(k\text{ even and }n=4k-2j+1\bigr).
\end{equation}
\end{conjecture}
The implication from \eqref{eq:only-zeros} to \eqref{eq:classification}
is proved above. The converse remains unproved in the following precise
region:
\begin{equation}\label{eq:gap-region}
 m=k-j\ge3,\qquad n\ge2k+2,
\end{equation}
away from the reflection zeros. All other regions have been handled by
the preceding analytic theorems.

\begin{theorem}[Exact reduction of the OEIS conjecture]\label{thm:equivalence}
The equality $a(n)=\U(n)$ for every $n\ge0$ is equivalent to
Conjecture~\ref{conj:support}. It is also equivalent to proving the
nonvanishing assertion only in the remaining region
\eqref{eq:gap-region}, excluding the reflection-zero family.
\end{theorem}
\begin{proof}
Formula~\eqref{eq:difference} identifies the expression in
\eqref{eq:classification} with the actual derivative coefficient.
Theorem~\ref{thm:upper} counts exactly the cells left after removal of the
two proved, disjoint zero families. Hence any additional zero makes the
corresponding inequality strict, and no additional zeros makes it an
equality. Theorem~\ref{thm:positive}, Proposition~\ref{prop:top}, and
Theorems~\ref{thm:m1}--\ref{thm:m2} eliminate the complementary regions.
\end{proof}

A tempting shortcut would assert that a derivative of a falling factorial
cannot vanish at an integer other than its reflection center. That
assertion is false. Direct differentiation gives
\begin{equation}\label{eq:false-shortcut}
 \frac{d^5}{du^5}\fall{u}{8}
 =3360(u-5)(u-2)(2u-7).
\end{equation}
It vanishes at $u=2$ and $u=5$, as well as at the reflection center
$u=7/2$. This is \emph{not} a counterexample to
Conjecture~\ref{conj:support}, because the derivative order, evaluation
point, and finite-difference order in that conjecture are linked in a
special way. It does show why those constraints must be used in a genuine
proof of the missing implication.

The transformed formula provides an alternative target. With
$d=n-2k-1\ge1$, the outstanding task is to prove that
\begin{equation}\label{eq:gap-transform}
 [u^{k+d+1}](1-u)^d(1+u)^{2m}
       \left(\frac{\atanh u}{u}\right)^m
\end{equation}
is nonzero for $m\ge3$, except when $d=2m$ and $k$ is even. The
reflection case is transparent because the series then becomes even;
the remaining finite-difference cancellations are the essential issue.
Equation~\eqref{eq:gap-transform} is an equivalent residual lemma, not an
additional theorem claimed by this article.

\section{Conclusion}
The proposed formula has a rigorous structural explanation. The obvious
degree loss removes the top logarithmic coefficients with $n>2k$.
A separate reflection symmetry removes exactly
$\lfloor(n+1)/8\rfloor$ further coefficients for odd $n$. Their total
produces the OEIS floor expression and its rational generating function
as an unconditional upper bound.

A fractional-linear transformation proves that a large region of the
coefficient triangle is strictly positive, and beta-integral arguments
exclude additional zeros on the first two deficit diagonals. These
results prove quadratic growth and reduce the full formula to the
explicit residual nonvanishing problem \eqref{eq:gap-transform}.
Exact finite computations settle every derivative order through $3000$.
The global equality is not proved here: its remaining obligation is the
nonvanishing statement in \eqref{eq:gap-region}, not any of the counting
or generating-function manipulations.

\appendix
\section{Reproducing the computations}\label{app:reproduce}
The archive needs Python 3.10 or later and a C++17 compiler for the
computations. Python uses only its standard library. The PDF can be
rebuilt with a conventional LaTeX installation containing the packages
listed in the source preamble.

Run the following commands from the archive directory. The Makefile
wraps the same operations.
\begin{lstlisting}
python3 code/verify_exact.py --max-n 200 --formula-n 16
mkdir -p build
g++ -O3 -std=c++17 -Wall -Wextra \
    code/verify_modular.cpp -o build/verify_modular
build/verify_modular 3000 1000000007 data/p1000000007 \
    data/watch_coordinates.csv
build/verify_modular 3000 1000000009 data/p1000000009 \
    data/watch_coordinates.csv
python3 code/combine_certificates.py
pdflatex -interaction=nonstopmode -halt-on-error \
    A290268_partial_results.tex
pdflatex -interaction=nonstopmode -halt-on-error \
    A290268_partial_results.tex
\end{lstlisting}
The watch file requests explicit residues at the seven exceptional
coordinates found in the modular runs. It does not restrict the search:
the C++ program checks every coefficient in every row regardless of the
watch file. The combiner checks that the complete exceptional zero lists
are disjoint and that all listed witnesses are consistent with them.

For a smaller test, the optional watch file can be omitted. For example,
\texttt{build/verify\_modular 100 1000000007 data/test} performs the full
modular recurrence through $100$. A watch file containing indices beyond
the chosen bound is rejected, rather than silently ignored.

\section{Files included in the archive}
\begin{longtable}{@{}>{\raggedright\arraybackslash}p{0.43\textwidth}>{\raggedright\arraybackslash}p{0.53\textwidth}@{}}
\toprule
File or directory & Purpose\\
\midrule
\endhead
\texttt{A290268\_partial\_results.tex} & Complete article source.\\
\texttt{A290268\_partial\_results.pdf} & Compiled article.\\
\texttt{README.md}, \texttt{Makefile} & Status, dependency notes, and reproduction commands.\\
\texttt{code/verify\_exact.py} & Arbitrary-precision recurrence and three independent coefficient-formula checks.\\
\texttt{code/verify\_modular.cpp} & Complete modular coefficient search with explicit overflow bounds.\\
\texttt{code/combine\_certificates.py} & Verification and combination of the two finite modular certificates.\\
\texttt{data/exact\_counts.csv} & Exact term counts through $200$.\\
\texttt{data/certified\_counts.csv} & Certified term counts through $3000$.\\
\texttt{data/p1000000007\_*} and \texttt{data/p1000000009\_*} & Complete exceptional zero lists, row statistics, witness residues, run logs, and summaries.\\
\texttt{data/modular\_witnesses.csv} & The seven coordinates at which one residue vanishes, with the nonzero residue in the other modulus.\\
\texttt{data/*summary.json} & Machine-readable verification scope and results.\\
\texttt{data/sample\_coefficients.json} & Fully collected low-order coefficient arrays and cancellation examples.\\
\texttt{SHA256SUMS.txt} & Integrity checks for the distributed files; these are not mathematical proof certificates.\\
\bottomrule
\end{longtable}

\begin{thebibliography}{9}
\bibitem{oeis}
The OEIS Foundation Inc.,
\emph{A290268: Number of terms in the fully expanded $n$-th derivative of
$x^{x^2}$}, contributed by Vladimir Reshetnikov, with formula comments by
Peter Luschny, October 2017. Accessed September 20, 2026.
\url{https://oeis.org/A290268}.

\bibitem{poon}
Steven S. Poon,
\emph{Higher Derivatives of the Falling Factorial and Related
Generalizations of the Stirling and Harmonic Numbers},
arXiv:1401.2737, 2014.
\url{https://arxiv.org/abs/1401.2737}.

\bibitem{dlmf-reflection}
NIST Digital Library of Mathematical Functions,
\S5.5(ii), especially Eq.~5.5.3, Euler's reflection formula.
Accessed September 20, 2026.
\url{https://dlmf.nist.gov/5.5.E3}.

\bibitem{dlmf-beta}
NIST Digital Library of Mathematical Functions,
\S5.12, especially Eq.~5.12.1, Euler's beta integral.
Accessed September 20, 2026.
\url{https://dlmf.nist.gov/5.12.E1}.
\end{thebibliography}
\end{document}
